\documentclass[12pt, titlepage]{article}

\usepackage{booktabs}
\usepackage{tabularx}
\usepackage{hyperref}
\hypersetup{
    colorlinks,
    citecolor=black,
    filecolor=black,
    linkcolor=red,
    urlcolor=blue
}
\usepackage[round]{natbib}

\input{../Comments}
\input{../Common}

\begin{document}

\title{Verification and Validation Report: \progname} 
\author{\authname}
\date{\today}
	
\maketitle

\pagenumbering{roman}

\section{Revision History}

\begin{tabularx}{\textwidth}{p{3cm}p{2cm}X}
\toprule {\bf Date} & {\bf Version} & {\bf Notes}\\
\midrule
Date 1 & 1.0 & Notes\\
Date 2 & 1.1 & Notes\\
\bottomrule
\end{tabularx}

~\newpage

\section{Symbols, Abbreviations and Acronyms}

\renewcommand{\arraystretch}{1.2}
\begin{tabular}{l l} 
  \toprule		
  \textbf{symbol} & \textbf{description}\\
  \midrule 
  T & Test\\
  \bottomrule
\end{tabular}\\

\wss{symbols, abbreviations or acronyms -- you can reference the SRS tables if needed}

\newpage

\tableofcontents

\listoftables %if appropriate

\listoffigures %if appropriate

\newpage

\pagenumbering{arabic}

This document ...

\section{Functional Requirements Evaluation}

\section{Nonfunctional Requirements Evaluation}

\subsection{Usability}
		
\subsection{Performance}

\subsection{etc.}
	
\section{Comparison to Existing Implementation}	

This section will not be appropriate for every project.

\section{Unit Testing}

\section{Changes Due to Testing}

\wss{This section should highlight how feedback from the users and from 
the supervisor (when one exists) shaped the final product.  In particular 
the feedback from the Rev 0 demo to the supervisor (or to potential users) 
should be highlighted.}

\section{Automated Testing}
		
\section{Trace to Requirements}

This section provides traceability between the requirements defined in the
Software Requirements Specification (SRS) and the verification activities
performed during testing. Each requirement is linked to one or more test
cases described in the Verification and Validation Plan. The results of the
tests demonstrate whether the implementation satisfies the specified
requirements.

Functional requirements are validated primarily through automated system
tests that simulate user interactions with the application. These tests
verify core functionality such as creating notes, manipulating diagram
elements, and performing file management operations.

Non-functional requirements are evaluated through a combination of manual
tests, usability evaluation, and performance measurements. These tests assess
system responsiveness, usability, cross-platform compatibility, and data
integrity.

Table~\ref{tab:reqTrace} summarizes the mapping between SRS requirements and
the test cases used to verify them.

\begin{table}[htbp]
\centering
\begin{tabular}{p{2cm} p{7cm} p{3cm} p{2cm}}
\toprule
\textbf{Req ID} & \textbf{Requirement Description} & \textbf{Test Case} & \textbf{Result} \\
\midrule

F211 & The system shall allow users to create and edit notes. 
& test-F211 & Pass \\

F214 & The system shall allow users to insert and manipulate shapes on the canvas. 
& test-F214 & Pass \\

F217 & The system shall allow users to save and reopen notes. 
& test-F217 & Pass \\

F243 & The system shall allow users to export notes as PDF files. 
& test-F243 & Pass \\

NF211 & The system shall provide low-latency input and rendering during editing operations. 
& test-NF211 & Pass \\

NF218 & The system shall maintain acceptable performance when handling a large number of shapes or objects. 
& test-NF218 & Pass \\

NF213 & The system shall support execution on Windows, macOS, and Linux platforms. 
& test-NF213 & Pass \\

NF214 & The system shall provide accessible keyboard shortcuts for editing operations. 
& test-NF214 & Pass \\

NF217 & The system shall provide an intuitive and consistent user interface. 
& test-NF217 & Pass \\

NF212 & The system shall maintain local data persistence and autosave functionality. 
& test-NF212 & Pass \\

NF215 & The system shall maintain file integrity during save and export operations. 
& test-NF215 & Pass \\

NF216 & The system shall not perform background network communication. 
& test-NF216 & Pass \\

\bottomrule
\end{tabular}
\caption{Traceability Between Requirements and Test Cases}
\label{tab:reqTrace}
\end{table}

The results indicate that all implemented requirements were successfully
verified through the corresponding tests. No critical failures were observed
during the verification process, demonstrating that the implemented system
satisfies the functional and non-functional requirements specified in the
SRS.

\vspace{25em}
		
\section{Trace to Modules}		

\section{Code Coverage Metrics}

\bibliographystyle{plainnat}
\bibliography{../../refs/References}

\newpage{}
\section*{Appendix --- Reflection}

The information in this section will be used to evaluate the team members on the
graduate attribute of Reflection.

\input{../Reflection.tex}

\begin{enumerate}
  \item What went well while writing this deliverable? 
  \item What pain points did you experience during this deliverable, and how
    did you resolve them?
  \item Which parts of this document stemmed from speaking to your client(s) or
  a proxy (e.g. your peers)? Which ones were not, and why?
  \item In what ways was the Verification and Validation (VnV) Plan different
  from the activities that were actually conducted for VnV?  If there were
  differences, what changes required the modification in the plan?  Why did
  these changes occur?  Would you be able to anticipate these changes in future
  projects?  If there weren't any differences, how was your team able to clearly
  predict a feasible amount of effort and the right tasks needed to build the
  evidence that demonstrates the required quality?  (It is expected that most
  teams will have had to deviate from their original VnV Plan.)
\end{enumerate}

\end{document}